% ═══════════════════════════════════════════════════════════════════════════
% ALTERNATIVE B: SPEED-DATING FORMAT
% Sicherer im Timing durch feste 3-Minuten-Intervalle
% ═══════════════════════════════════════════════════════════════════════════

\documentclass[11pt, a4paper]{article}

\usepackage[utf8]{inputenc}
\usepackage[T1]{fontenc}
\usepackage[ngerman]{babel}
\usepackage{helvet}
\renewcommand{\familydefault}{\sfdefault}

\usepackage[margin=2.5cm, left=3.5cm]{geometry}
\usepackage{setspace}
\onehalfspacing

\usepackage{enumitem}
\usepackage{tabularx}
\usepackage{booktabs}
\usepackage{array}
\usepackage{longtable}

\usepackage[table]{xcolor}
\usepackage{tcolorbox}
\usepackage{fancyhdr}
\usepackage{lastpage}
\usepackage{titlesec}
\usepackage{parskip}
\usepackage{tikz}
\usepackage{amssymb}
\usepackage[hidelinks]{hyperref}

% Farben
\definecolor{spanisch}{HTML}{c41e3a}
\definecolor{grau}{HTML}{666666}
\definecolor{hellgrau}{HTML}{f5f5f5}
\definecolor{success}{HTML}{2a7a4b}
\definecolor{warning}{HTML}{b35c00}
\definecolor{alternativeB}{HTML}{1565c0}

\titleformat{\section}{\normalfont\large\bfseries\color{spanisch}}{\thesection}{1em}{}
\titleformat{\subsection}{\normalfont\normalsize\bfseries}{\thesubsection}{1em}{}

\pagestyle{fancy}
\fancyhf{}
\lhead{\small Unterrichtslangentwurf -- Spanisch Spätbeginner}
\rhead{\small DS 10 -- \textcolor{alternativeB}{Alternative B: Speed-Dating}}
\cfoot{\small Seite \thepage\ von \pageref{LastPage}}

\newcommand{\es}[1]{\textcolor{spanisch}{\textit{#1}}}

\newtcolorbox{infobox}[1][]{colback=hellgrau, colframe=spanisch, boxrule=0.8pt, arc=0pt, left=8pt, right=8pt, top=6pt, bottom=6pt, #1}

\newtcolorbox{alternativebox}[1][]{colback=alternativeB!10, colframe=alternativeB, boxrule=1.5pt, arc=0pt, left=8pt, right=8pt, top=6pt, bottom=6pt, title={\textbf{ALTERNATIVE B: SPEED-DATING}}, fonttitle=\color{alternativeB}\bfseries, #1}

\begin{document}

% ═══════════════════════════════════════════════════════════════════════════
% DECKBLATT
% ═══════════════════════════════════════════════════════════════════════════

\thispagestyle{empty}

\begin{center}
    \vspace*{1cm}

    {\Large\bfseries Unterrichtslangentwurf}

    \vspace{0.3cm}

    {\large Tier 1 -- Vollständig}

    \vspace{0.5cm}

    \fcolorbox{alternativeB}{alternativeB!10}{%
        \parbox{0.6\textwidth}{%
            \centering\vspace{0.3cm}
            {\large\color{alternativeB}\textbf{ALTERNATIVE B}}\\[0.2cm]
            {\Large\bfseries Speed-Dating: ¿Quién eres?}\\[0.2cm]
            {\normalsize Strukturierte Kurzinterviews mit Partnerwechsel}
            \vspace{0.3cm}
        }%
    }

    \vspace{1cm}

    \begin{tabular}{ll}
        \textbf{Lehrkraft:} & [Name] \\[6pt]
        \textbf{Schule:} &  \\[6pt]
        \textbf{Kurs:} & Spanisch Spätbeginner (Kl. 10-12) \\[6pt]
        \textbf{Datum:} & Dienstag, 03.02.2026 \\[6pt]
    \end{tabular}

    \vspace{1.5cm}

    \fcolorbox{success}{white}{%
        \parbox{0.8\textwidth}{%
            \centering\vspace{0.3cm}
            {\large\color{success}\textbf{VORTEILE DIESER VARIANTE}}\\[0.3cm]
            \begin{tabular}{ll}
                $\checkmark$ & Feste 3-Min-Intervalle = kein Zeitdruck \\
                $\checkmark$ & Alle SuS sprechen mit mehreren Partnern \\
                $\checkmark$ & Kein Notieren während des Gesprächs \\
                $\checkmark$ & Einfache Struktur, wenig kann schiefgehen \\
            \end{tabular}
            \vspace{0.3cm}
        }%
    }

    \vfill

    {\small\color{grau} Risikobewertung: \textbf{NIEDRIG} -- Timing sehr sicher}

\end{center}

\newpage

\tableofcontents
\newpage

% ═══════════════════════════════════════════════════════════════════════════
% KONZEPT
% ═══════════════════════════════════════════════════════════════════════════

\section{Konzept: Speed-Dating}

\begin{alternativebox}
\textbf{Grundidee:} Anstatt langer Tandem-Interviews mit Notieren und Präsentation nutzt diese Variante das Speed-Dating-Format: Kurze, strukturierte Gespräche (3 Min) mit festem Partnerwechsel.

\textbf{Struktur:}
\begin{enumerate}[nosep]
    \item \textbf{3 Minuten} Gespräch mit Partner
    \item \textbf{Gong/Signal} → Wechsel zum nächsten Partner
    \item Nach 5-6 Runden: Reflexion im Plenum
\end{enumerate}

\textbf{Vorteile für Beobachtung:}
\begin{itemize}[nosep]
    \item Hohe Sprechaktivität (alle sprechen gleichzeitig)
    \item Klare Struktur (Timer steuert alles)
    \item Keine Schreibphasen = kein Zeitrisiko durch langsame Schreiber
    \item Natürliche Differenzierung (A spricht mit B, B unterstützt)
\end{itemize}
\end{alternativebox}

\subsection{Warum kein Notieren?}

Das Notieren während des Gesprächs ist die Hauptursache für Zeitverzögerungen:
\begin{itemize}
    \item Langsame Schreiber bremsen das Gespräch
    \item Fokus liegt auf Schreiben statt Sprechen
    \item Gespräch wirkt künstlich („Warte, ich schreibe noch...")
\end{itemize}

\textbf{Lösung:} SuS merken sich 2-3 Fakten pro Partner, notieren NACH dem Speed-Dating auf einer Sammelkarte.

\subsection{Sitzordnung}

\begin{center}
\begin{tikzpicture}
\draw[thick] (0,0) -- (10,0);
\foreach \x in {1,3,5,7,9} {
    \fill[spanisch] (\x,0.3) circle (0.3);
    \node[white, font=\small\bfseries] at (\x,0.3) {A};
}
\foreach \x in {1,3,5,7,9} {
    \fill[alternativeB] (\x,-0.3) circle (0.3);
    \node[white, font=\small\bfseries] at (\x,-0.3) {B};
}
\draw[->, thick] (9.5,-0.3) arc (-90:90:0.6);
\node[right] at (10.2,0) {B rückt weiter};
\end{tikzpicture}
\end{center}

\textbf{Aufbau:} Zwei Reihen gegenüber. Nach jedem Signal rückt die B-Reihe einen Platz weiter.

% ═══════════════════════════════════════════════════════════════════════════
% STUNDENZIELE
% ═══════════════════════════════════════════════════════════════════════════

\section{Stundenziele}

\subsection{Grobziel}

Die SuS erweitern ihre \textbf{mündliche Interaktionskompetenz}, indem sie in strukturierten Kurzgesprächen Informationen über sich selbst austauschen und Fragen an verschiedene Partner stellen.

\subsection{Feinziele}

\begin{center}
\begin{tabularx}{\textwidth}{|c|X|c|c|}
    \hline
    \textbf{Nr.} & \textbf{Die SuS...} & \textbf{AFB} & \textbf{Diff.} \\
    \hline
    FZ1 & ...stellen mindestens 3 vorgegebene Interview-Fragen korrekt. & I & alle \\
    \hline
    FZ2 & ...beantworten Fragen zu ihrer Person flüssig und korrekt. & II & alle \\
    \hline
    FZ3 & ...passen ihre Sprechgeschwindigkeit an den Partner an. & II & B \\
    \hline
    FZ4 & ...erinnern sich an 2-3 Fakten pro Partner für die Reflexion. & II & alle \\
    \hline
    FZ5 & ...berichten über einen Partner im Plenum (3. Person). & III & B \\
    \hline
\end{tabularx}
\end{center}

% ═══════════════════════════════════════════════════════════════════════════
% VERLAUFSPLANUNG
% ═══════════════════════════════════════════════════════════════════════════

\section{Verlaufsplanung (Beobachtete 45 Min)}

\subsection{Übersicht}

\begin{center}
\begin{tabular}{|l|c|l|}
    \hline
    \textbf{Phase} & \textbf{Zeit} & \textbf{Aktivität} \\
    \hline
    Orientierung & 5 Min & Regeln, Sitzordnung \\
    \hline
    Speed-Dating & 24 Min & 6 Runden à 3 Min + Wechsel \\
    \hline
    Notieren & 5 Min & Sammelkarte ausfüllen \\
    \hline
    Reflexion & 8 Min & Berichte im Plenum \\
    \hline
    Cierre & 3 Min & Abschluss \\
    \hline
    \textbf{Total} & \textbf{45 Min} & \\
    \hline
\end{tabular}
\end{center}

\subsection{Detaillierte Verlaufsplanung}

\begin{longtable}{|p{1.2cm}|p{2cm}|p{4cm}|p{3.5cm}|p{1.5cm}|}
    \hline
    \textbf{Zeit} & \textbf{Phase} & \textbf{Lehrerhandeln} & \textbf{Schülerhandeln} & \textbf{SF} \\
    \hline
    \endfirsthead

    0--5 & Orien-tierung & Begrüßung. Erklärt Speed-Dating: „3 Minuten Gespräch, dann Gong und Wechsel. Kein Schreiben während des Gesprächs!" Zeigt Sitzordnung. Verteilt Fragekarten. & Hören zu, bilden zwei Reihen (A gegenüber B) & PL \\
    \hline

    \multicolumn{5}{|c|}{\cellcolor{hellgrau}\textbf{SPEED-DATING (6 Runden × 3 Min = 18 Min + 6 Min Wechsel)}} \\
    \hline

    5--8 & Runde 1 & Timer 3 Min. Geht herum, beobachtet. & A und B stellen sich gegenseitig Fragen & PA \\
    \hline

    8--9 & Wechsel & Gong. „¡Cambio!" B-Reihe rückt einen Platz weiter. & B wechselt zum nächsten A & -- \\
    \hline

    9--12 & Runde 2 & Timer 3 Min. & Neue Paare sprechen & PA \\
    \hline

    12--13 & Wechsel & Gong. „¡Cambio!" & B wechselt & -- \\
    \hline

    13--16 & Runde 3 & Timer 3 Min. & Neue Paare sprechen & PA \\
    \hline

    16--17 & Wechsel & Gong. „¡Cambio!" & B wechselt & -- \\
    \hline

    17--20 & Runde 4 & Timer 3 Min. & Neue Paare sprechen & PA \\
    \hline

    20--21 & Wechsel & Gong. „¡Cambio!" & B wechselt & -- \\
    \hline

    21--24 & Runde 5 & Timer 3 Min. & Neue Paare sprechen & PA \\
    \hline

    24--25 & Wechsel & Gong. „¡Cambio!" & B wechselt & -- \\
    \hline

    25--28 & Runde 6 & Timer 3 Min. Letzte Runde! & Neue Paare sprechen & PA \\
    \hline

    28--29 & Ende & „¡Stop! Muy bien. Ahora a los asientos." & Zurück zu normalen Plätzen & -- \\
    \hline

    \multicolumn{5}{|c|}{\cellcolor{hellgrau}\textbf{NOTIEREN (5 Min)}} \\
    \hline

    29--34 & Sammel-karte & Verteilt Sammelkarte. „Escribe 2-3 cosas que recuerdas de cada compañero." Timer 5 Min. & Füllen Sammelkarte aus (Namen + 2-3 Fakten) & EA \\
    \hline

    \multicolumn{5}{|c|}{\cellcolor{hellgrau}\textbf{REFLEXION (8 Min)}} \\
    \hline

    34--42 & Plenum & Fragt: \es{¿Qué has aprendido sobre...?} Ruft 4-5 SuS auf. Hilft bei 3. Person. & B berichtet über A-Partner (3. Person): \es{Se llama..., es..., le gusta...} & PL \\
    \hline

    \multicolumn{5}{|c|}{\cellcolor{hellgrau}\textbf{ABSCHLUSS (3 Min)}} \\
    \hline

    42--45 & Cierre & Zusammenfassung: \es{¡Muy bien! Ahora conocemos mejor a nuestros compañeros.} Feedback, Verabschiedung. & Hören zu & PL \\
    \hline

\end{longtable}

% ═══════════════════════════════════════════════════════════════════════════
% MATERIALIEN
% ═══════════════════════════════════════════════════════════════════════════

\section{Materialien}

\subsection{Fragekarten (vereinfacht)}

\textbf{Für alle SuS dieselbe Karte} (keine A/B-Differenzierung bei den Fragen):

\begin{infobox}
\textbf{Speed-Dating Fragen:}

\begin{enumerate}
    \item \es{¿Cómo te llamas?}
    \item \es{¿Cómo eres?} (simpático, creativo, inteligente...)
    \item \es{¿Qué te gusta hacer?}
    \item \es{¿Qué profesión quieres tener?} (opcional für Fortgeschrittene)
\end{enumerate}

\textbf{Hilfe:} \es{Me llamo... / Soy... / Me gusta... / Quiero ser...}
\end{infobox}

\textbf{Differenzierung:} A-SuS nutzen Fragen 1-3, B-SuS alle 4 Fragen.

\subsection{Sammelkarte (nach dem Speed-Dating)}

\begin{center}
\begin{tabular}{|p{3cm}|p{4cm}|p{4cm}|}
    \hline
    \textbf{Nombre} & \textbf{Es...} & \textbf{Le gusta...} \\
    \hline
    & & \\[0.8cm]
    \hline
    & & \\[0.8cm]
    \hline
    & & \\[0.8cm]
    \hline
    & & \\[0.8cm]
    \hline
    & & \\[0.8cm]
    \hline
    & & \\[0.8cm]
    \hline
\end{tabular}
\end{center}

% ═══════════════════════════════════════════════════════════════════════════
% RISIKOBEWERTUNG
% ═══════════════════════════════════════════════════════════════════════════

\section{Risikobewertung}

\begin{center}
\begin{tabular}{|l|c|p{6cm}|}
    \hline
    \textbf{Risiko} & \textbf{Level} & \textbf{Begründung} \\
    \hline
    Timing & \textcolor{success}{\textbf{NIEDRIG}} & Feste 3-Min-Intervalle, Timer steuert alles \\
    \hline
    Schreibgeschwindigkeit & \textcolor{success}{\textbf{KEINS}} & Kein Schreiben während Gespräch \\
    \hline
    Langsame Gespräche & \textcolor{success}{\textbf{NIEDRIG}} & 3 Min pro Gespräch ist ausreichend \\
    \hline
    Präsentationsdauer & \textcolor{success}{\textbf{NIEDRIG}} & Nur kurze Berichte (1-2 Sätze pro SuS) \\
    \hline
    Nervosität & \textcolor{warning}{\textbf{MITTEL}} & Viele Partnerwechsel = weniger Tiefe \\
    \hline
\end{tabular}
\end{center}

\subsection{Gesamtrisiko}

\begin{center}
\fcolorbox{success}{success!10}{%
\parbox{0.8\textwidth}{%
\centering\vspace{0.3cm}
{\large\color{success}\textbf{GESAMTRISIKO: NIEDRIG}}\\[0.3cm]
Diese Variante ist timing-sicher. Der Timer steuert den gesamten Ablauf.\\
Keine Phase ist von individueller Schreibgeschwindigkeit abhängig.
\vspace{0.3cm}
}%
}
\end{center}

% ═══════════════════════════════════════════════════════════════════════════
% DIDAKTIK-SCORE
% ═══════════════════════════════════════════════════════════════════════════

\section{Didaktik-Score}

\begin{center}
\begin{tabular}{|c|l|c|p{5cm}|}
    \hline
    \textbf{\#} & \textbf{Dimension} & \textbf{Score} & \textbf{Notiz} \\
    \hline
    1 & Differenzierung & 8/10 & Implizit durch Fragetiefe, weniger explizit als Variante A \\
    \hline
    2 & Taxonomie (AFB) & 8/10 & Hauptsächlich AFB I-II, weniger AFB III \\
    \hline
    3 & Lernziel-Alignment & 9/10 & Kommunikative Kompetenz = Kernziel \\
    \hline
    4 & Aktivierung & 10/10 & Alle sprechen in jeder Runde \\
    \hline
    5 & Scaffolding & 8/10 & Fragekarten als Hilfe, aber weniger gestuft \\
    \hline
    6 & Feedback-Qualität & 7/10 & Weniger individuelles Feedback während Runden \\
    \hline
    7 & Sprachliche Klarheit & 10/10 & Sehr klare, einfache Struktur \\
    \hline
    8 & Motivation & 9/10 & Speed-Dating-Format ist motivierend \\
    \hline
    9 & Inklusion (UDL) & 8/10 & Einheitliche Struktur für alle \\
    \hline
    10 & Praktikabilität & 10/10 & Timing sehr sicher, wenig Vorbereitung \\
    \hline
    \multicolumn{2}{|r|}{\textbf{GESAMT:}} & \textbf{8.7/10} & \textcolor{success}{\textbf{FREIGABE}} \\
    \hline
\end{tabular}
\end{center}

\vspace{0.5cm}

\textbf{Fazit:} Diese Variante ist \textbf{timing-sicherer} als die Original-Variante, zeigt aber weniger explizite Differenzierung. Geeignet, wenn Zeitsicherheit oberste Priorität hat.

\end{document}
